\documentclass[11pt]{article}

% 中文注释:
% - 模板沿用 ACL 提供的 acl.sty
% - review 模式投稿;最终版改成 \usepackage{acl}
% - 6 页正文目标(EMNLP demo / system 论文长版),refs 不计入
\usepackage[review]{acl}

\usepackage{times}
\usepackage{latexsym}
\usepackage[T1]{fontenc}
\usepackage[utf8]{inputenc}
\usepackage{microtype}
\usepackage{inconsolata}
\usepackage{graphicx}
\usepackage{booktabs}
\usepackage{xcolor}
\usepackage{listings}

% 中文注释: 代码块样式 - system prompt / JSON schema 展示用
\lstdefinestyle{jsonstyle}{
  basicstyle=\ttfamily\footnotesize,
  breaklines=true,
  columns=fullflexible,
  keepspaces=true,
  showstringspaces=false,
  frame=single,
  framerule=0.3pt,
  framesep=3pt,
}

% 中文注释: 标题候选 3 个,A 最有冲击力(hallucination-free 是钩子)
% A. GroundedGeo: Hallucination-Free Tool-Use for Mineral Prospectivity Reasoning
% B. Anchoring LLMs in Expert Trees: A Grounded Agent for Mineral Exploration
% C. From Pathfinders to Prompts: A Faithful NL Interface for Multi-Commodity Prospectivity Mapping
\title{GroundedGeo: Hallucination-Free Tool-Use \\
       for Multi-Commodity Mineral Prospectivity Reasoning}

% 中文注释: 投稿时 anonymous,camera-ready 再改回真实作者
\author{Anonymous Submission \\
  \texttt{anonymous@emnlp.org} \\}

\begin{document}
\maketitle

% =====================================================================
% Abstract
% =====================================================================
\begin{abstract}
% 中文注释: Abstract 必须命中 4 个点(每个 1-2 句话):
%   1. 痛点: LLM 在数值密集科学领域幻觉机理 / 数字
%   2. 任务: 多金属矿产潜力推理 - 数字 + 机理双重可证伪
%   3. 方法: grounded tool-use, 100 case study 库, 9 金属 x 9 区域
%   4. 结果: hallucination rate 降 X%, 人评偏好显著
% 中文注释: 不要太学术腔,审稿人看 abstract 30 秒就要决定要不要点开
LLM agents increasingly assist scientific workflows, but in quantitative
domains such as mineral exploration they routinely hallucinate
measurements, mechanisms, and citations that domain experts cannot verify
against the underlying data. We present \textsc{GroundedGeo}, a tool-use
agent for multi-commodity mineral prospectivity reasoning over Western
Australia, in which every numeric and explanatory claim must originate
from a structured tool return rather than free-form generation.
The system wraps an existing expert-tree scoring pipeline across nine
commodities (Cu, Au, Ni, W, Sn, Co, Ta, Mn, REE) with ten typed tools
spanning point queries, regional reconnaissance, and analog search.
We pre-compute a library of approximately one hundred
(region~$\times$~commodity) case studies covering five major cratons and
four mining regions, enabling latency-free named-region queries while
supporting custom bounding-box and circular regions through live
two-stage grid scoring. In a faithfulness evaluation over 30
expert-authored prompts, \textsc{GroundedGeo} reduces numeric
hallucinations by \textbf{XX\%} relative to prompt-only GPT-4o while
preserving explanatory fluency, and is preferred by geology students in
pairwise comparison ($p < 0.01$). The system and case study library are
released at \url{https://anonymous.4open.science/r/groundedgeo}.
\end{abstract}

% =====================================================================
% §1 Introduction (~3/4 page)
% =====================================================================
\section{Introduction}
\label{sec:intro}

% 中文注释: Intro 第一段 - "LLM agent 已经无处不在,在 X / Y / Z 领域很成功"
% 中文注释: 引文: ReAct, Toolformer, AutoGPT, function-calling APIs
LLM-based agents now plan and orchestrate tools across software
engineering, web automation, and general knowledge tasks
\citep{yao2022react,schick2023toolformer,openai2023functioncalling}.
A common pattern is to expose an LLM to a set of typed function
signatures and let it decide \emph{which} tool to call \emph{when},
yielding chains of structured operations stitched together with natural
language.

% 中文注释: 第二段 - "但是在 quantitative scientific domain 这种 pattern 出问题"
% 中文注释: 关键论点: 每个数字都可证伪、每个机理都有学界共识 -- 编不得
However, when the underlying task is \emph{quantitatively grounded}---
when every numeric claim is verifiable against measured data and every
proposed mechanism is subject to established scientific consensus---
prompt-only LLMs systematically hallucinate.
This failure mode is well documented in medical
\citep{singhal2023large}, legal \citep{dahl2024large}, and scientific
\citep{petroni2021kilt} domains, where post-hoc verification is hard
because the LLM's output is interleaved natural language rather than
discrete tool calls.

% 中文注释: 第三段 - 把焦点拉到 mineral exploration,这是个 sharp test case
% 中文注释: Figure 1 在这里出现(motivating example)
Mineral exploration is an unusually sharp test bed for grounded
scientific reasoning. Every prospectivity assessment can be broken down
into measurable feature anomalies (geochemical z-scores, magnetic lows,
fault proximity) and a small set of well-established mineralisation
models. A claim that a region is prospective for tungsten because
``W/Mo ratios are anomalously high alongside Bi enrichment'' is either
true in the data or it is not. Figure~\ref{fig:bad-llm} contrasts a
typical prompt-only GPT-4o response to a Western Australia coordinate
with the actual measured signals at that point: the model fabricates
specific numerical anomalies, invokes pathfinder elements absent from
the local sampling, and asserts geological context not supported by the
regional dataset.

\begin{figure}[t]
  % 中文注释: Figure 1 用 'bad LLM 反例'
  % 中文注释: 上半: 裸 GPT-4o 输出, 红色高亮虚构数字 + 虚构 pathfinder
  % 中文注释: 下半: 同点真实 sediment assay + active experts list
  % 中文注释: 必须有视觉冲击力 - 这是审稿人决定继续往下读的钩子
  \centering
  \includegraphics[width=\columnwidth]{example-image-golden}
  \caption{\textbf{Motivating example.} Top: a prompt-only GPT-4o
  response when asked to assess W (tungsten) prospectivity at
  $(121.4^\circ\text{E}, 28.3^\circ\text{S})$. Red highlights mark
  numerical and mechanistic claims absent from or inconsistent with the
  underlying WAMEX dataset. Bottom: the same point's actual top
  feature signals retrieved via our grounded agent.}
  \label{fig:bad-llm}
\end{figure}

% 中文注释: 第四段 - 我们的方案 1 句话先说
% 中文注释: 然后过渡到 3 个 contributions
In this work we propose a tool-use architecture that
\emph{structurally} prevents these failures by funneling all
quantitative and mechanistic claims through a typed tool layer
backed by an expert-tree scoring pipeline. Rather than asking the LLM
to recall geological facts, we restrict it to invoking, composing, and
verbalising the returns of ten purpose-built tools.

% 中文注释: 3 个 contributions - 写清楚 NLP 含量
% 中文注释: 审稿人最容易拒的理由是 "this is a domain demo not NLP" -- 这里要正面应对
We make three contributions:

\begin{itemize}
  % 中文注释: 贡献 1 - 方法学
  \item A grounded tool-use architecture (Section~\ref{sec:architecture})
  that combines typed tools, audit-trail returns, and a faithfulness
  system prompt to make numeric and mechanistic hallucinations
  detectable and (in practice) absent.
  % 中文注释: 贡献 2 - 数据资源
  \item A library of approximately one hundred precomputed
  (region~$\times$~commodity) case studies
  (Section~\ref{sec:casestudy}) covering nine commodities across nine
  metallogenically distinct regions of Western Australia, enabling
  latency-free demo at scale.
  % 中文注释: 贡献 3 - 评估
  \item A faithfulness evaluation protocol (Section~\ref{sec:eval})
  that quantifies numeric hallucination, mechanism plausibility,
  and tool-call accuracy on 30 expert-authored prompts, combining
  automatic checks with pairwise human judgements.
\end{itemize}

% =====================================================================
% §2 Related Work (~1/4 page)
% =====================================================================
\section{Related Work}
\label{sec:related}

% 中文注释: 三段 - 每段一个领域,最后一段说 gap
% 中文注释: 不要写成 laundry list,要写出 gap

\paragraph{Tool-use and agentic LLMs.}
% 中文注释: ReAct, Toolformer, OpenAI function calling, Anthropic tool use
% 中文注释: Mind2Web, SWE-bench 这些 agent benchmark
ReAct \citep{yao2022react} interleaves reasoning traces with tool
invocations; Toolformer \citep{schick2023toolformer} learns when to
call tools via self-supervision; function-calling APIs from major
providers \citep{openai2023functioncalling} now make typed tool use a
first-class primitive. Recent agentic benchmarks
\citep{mialon2023gaia,liu2023agentbench} measure planning and
composition. None of these works specifically address
\emph{faithfulness-first} agent design in a quantitative scientific
domain.

\paragraph{Faithful and grounded generation.}
% 中文注释: Reiter & Dale 经典 NLG pipeline (重要,框架根基)
% 中文注释: 现代 grounding: retrieval-augmented, attribution, citation models
% 中文注释: hallucination evaluation: TruthfulQA, FActScore
Classical NLG pipelines \citep{reiter2000building} separate content
selection, structuring, and surface realisation, making each step
auditable. Recent grounded generation work attaches generated text to
retrieved passages \citep{lewis2020retrieval,bohnet2022attributed} or
verifies factual claims post-hoc \citep{min2023factscore}. Our work
draws on the classical pipeline philosophy but moves grounding
\emph{into} the agent loop: the LLM cannot generate ungrounded numeric
content because the relevant numbers are only available through tool
returns.

\paragraph{Machine learning for mineral prospectivity.}
% 中文注释: Carranza WofE, Zuo deep learning, recent SHAP-on-RF
% 中文注释: 强调他们的可解释性停留在 layer-level / feature attribution,不到 mechanism-level
Knowledge-driven prospectivity modelling has a long history
\citep{carranza2017improved}; data-driven approaches now include
random forests, gradient boosting, and deep CNNs
\citep{zuo2020deep,sun2020data}. Recent work attempts post-hoc
interpretation via SHAP \citep{wang2022deep}, but explanations remain
at the level of feature attribution rather than mechanistic geological
reasoning, and none expose a natural-language interface to domain
users.

% 中文注释: gap 段 - 三个领域的交叉点没人做
The intersection of these three threads---faithful generation at
agent-loop time, in a quantitatively grounded scientific domain, with
a natural-language interface for non-engineering users---has not been
systematically explored. \textsc{GroundedGeo} occupies this gap.

% =====================================================================
% §3 System Architecture (~1 page)
% =====================================================================
\section{System Architecture}
\label{sec:architecture}

% 中文注释: Figure 2 在这里 - 5 层架构图(从 5 层框架那张图改成 LaTeX 矢量)
\begin{figure*}[t]
  \centering
  \includegraphics[width=0.92\textwidth]{example-image-a}
  \caption{\textbf{\textsc{GroundedGeo} architecture.} A five-layer
  stack with explicit dual-path routing at Layer~2: named-region
  queries are served from a precomputed case-study library, while
  custom (bounding-box / circle) regions trigger live two-stage grid
  scoring. The LLM (Layer~3) interacts exclusively with structured
  JSON tool returns from Layer~2; it never reads Layer~0--1 directly.}
  \label{fig:arch}
\end{figure*}

% 中文注释: Layer 0/1 - 已有部分, 一段带过
\paragraph{Layers 0--1: Data and scoring engine.}
The lower layers are inherited from a prior multi-commodity
prospectivity pipeline. Layer~0 indexes WAMEX sediment and rockchip
assays, gravity and magnetic rasters, and structural vector layers
(faults, lineaments, geological provinces). Layer~1 instantiates one
\texttt{ProspectivityModel} per commodity, which fits a tree of typed
experts---geochemical pathfinders, enrichment, ratio anomalies,
magnetic lows, fault proximity---and produces a \texttt{NodeScore}
tree at inference time. The system supports nine commodities (Cu, Au,
Ni, W, Sn, Co, Ta, Mn, REE) with four negative-sampling strategies
validated against held-out known deposits; AUC scores under the
\emph{far-random} strategy range from 0.81 to 0.99 across commodities
(Table~\ref{tab:auc-summary}).

% 中文注释: AUC summary table (展示 9 金属的实证厚度)
\begin{table}[t]
  \centering
  \footnotesize
  \begin{tabular}{lcccc}
    \toprule
    \textbf{Metal} & \textbf{Rand.} & \textbf{Far-R.} & \textbf{NonMine} & \textbf{Spatial} \\
    \midrule
    Cu  & 0.861 & 0.925 & 0.629 & 0.576 \\
    Au  & 0.925 & 0.958 & 0.720 & 0.730 \\
    Ni  & 0.990 & 0.994 & 0.946 & 0.917 \\
    W   & 0.907 & 0.936 & 0.740 & 0.741 \\
    Sn  & 0.972 & 0.989 & 0.895 & 0.664 \\
    Co  & 0.977 & 0.983 & 0.929 & 0.805 \\
    Ta  & 0.919 & 0.953 & 0.775 & 0.782 \\
    Mn  & 0.802 & 0.825 & 0.745 & 0.257 \\
    REE & 0.814 & 0.807 & 0.795 & 0.577 \\
    \bottomrule
  \end{tabular}
  \caption{Scoring-engine AUC across nine commodities and four
  negative-sampling strategies (Section~\ref{sec:architecture}).
  Far-random keeps held-out positives but excludes negatives within
  50\,km of any known site.}
  \label{tab:auc-summary}
\end{table}

% 中文注释: Layer 2 - reasoning,narrative.describe 在这里
\paragraph{Layer 2: Structured reasoning.}
The \texttt{narrative.describe} module deterministically converts a
\texttt{NodeScore} tree to a \texttt{PointNarrative} object containing
(i) a global score and tier label, (ii) ranked
\texttt{ExpertContrib} entries, and (iii) ranked \texttt{FeatureSignal}
entries with z-scores, source attributions, and contributing-expert
identifiers. Crucially, \texttt{PointNarrative} is a structured
record, not a string: surface verbalisation is deferred to the agent
layer, which means the same record can be queried, compared, or
aggregated without re-running scoring.

% 中文注释: Layer 3 - 重点,agent layer
\paragraph{Layer 3: Grounded agent.}
The agent layer is the focus of this paper. It comprises four
modules:
(a)~\textbf{tools.py}, exposing ten typed functions
(Section~\ref{sec:tools}, Table~\ref{tab:tools});
(b)~\textbf{model\_cache.py}, lazily instantiating up to nine
\texttt{ProspectivityModel}s under LRU caching;
(c)~\textbf{runner.py}, implementing the OpenAI function-calling loop
with explicit audit-trail capture;
and (d)~\textbf{system\_prompt.py}, encoding faithfulness constraints
(Section~\ref{sec:prompt}). The LLM observes only the JSON tool
returns; raw assay tables, raster pixels, and vector geometries never
enter its context.

% 中文注释: Layer 4 - demo UI
\paragraph{Layer 4: Demonstration interface.}
A Gradio web application renders the agent loop as a three-pane view:
(left) chat history; (centre) an interactive Folium map with
clickable known sites, score heatmaps, and bounding-box / circle
selectors; (right) a collapsible tool-call trace showing each
invocation's arguments, return JSON, and audit fields. Screenshots
appear in Figure~\ref{fig:demo-ui}.

% =====================================================================
% §4 Grounded Tool Design (~3/4 page)
% =====================================================================
\section{Grounded Tool Design}
\label{sec:tools}

% 中文注释: 这一节是论文的方法学卖点 - 必须把工具设计的哲学讲清楚

% 中文注释: 三条设计原则
We adopt three design principles, each chosen to make hallucination
either structurally impossible or post-hoc detectable.

\paragraph{P1: Audit-trail returns.}
Every tool returns a JSON object containing an explicit
\texttt{audit} field that enumerates (i) the active data sources
loaded for this query, (ii) the active expert modules in the scored
tree, and (iii) the per-feature z-scores that contributed to the
returned score. Any quantitative claim downstream can be
mechanically traced to a specific field in some audit record.

\paragraph{P2: No free-form narrative slots.}
Only two tools (\texttt{explain\_point} and \texttt{explain\_region})
return text, and even then the text is structured: a sequence of
typed \texttt{FeatureSignal} records, each carrying its z-score and
source. The agent verbalises by selecting and re-ordering these
records; it cannot insert new signals.

\paragraph{P3: Numeric quotation requirement.}
The system prompt (Section~\ref{sec:prompt}) instructs the agent that
every numeric or pathfinder claim in its response must appear
verbatim in some tool return from the current turn. We verify this
constraint by post-hoc regex extraction on rendered responses
(Section~\ref{sec:eval}).

% 中文注释: Table 1 - 10 个工具列表
\begin{table*}[t]
  \centering
  \footnotesize
  \begin{tabular}{p{3.4cm}p{4.6cm}p{5.5cm}}
    \toprule
    \textbf{Tool} & \textbf{Inputs} & \textbf{Returns (key fields)} \\
    \midrule
    \texttt{list\_targets} & --- & commodity list, AUC summary \\
    \texttt{describe\_target} & \texttt{target} & pathfinders, ratio features, active layers \\
    \texttt{score\_point} & \texttt{target, lon, lat} & g-score, tier, audit \\
    \texttt{explain\_point} & \texttt{target, lon, lat, top\_n} & top experts, top signals, tier reason \\
    \texttt{score\_with\_direct\_evidence} & \texttt{target, lon, lat, assays} & combined score, evidence expert breakdown \\
    \texttt{multi\_target\_scan} & \texttt{lon, lat, targets} & ranked metals at point \\
    \texttt{find\_neighbors} & \texttt{lon, lat, radius\_km, kind} & nearby known sites or samples \\
    \texttt{resolve\_region} & named / bbox / circle spec & \texttt{Region} object with geometry \\
    \texttt{rank\_targets\_in\_region} & \texttt{region, top\_k} & ranked metals over region \\
    \texttt{scan\_region} & \texttt{region, target} & grid scores, top anomalies \\
    \texttt{top\_anomalies} & \texttt{region, target, top\_k, exclude\_known} & ranked points with explanations \\
    \texttt{explain\_region} & \texttt{region, target} & dominant pathfinders, regional narrative \\
    \bottomrule
  \end{tabular}
  \caption{The twelve grounded tools exposed to the agent. Audit
  fields are returned by every tool but omitted from this table for
  space.}
  \label{tab:tools}
\end{table*}

% 中文注释: 展开一个 example schema - explain_point
\paragraph{Example schema: \texttt{explain\_point}.}
We illustrate the design with the most-invoked tool.
% 中文注释: 这里贴一个简化的 JSON schema, 不是完整 OpenAI 那种带 description 的版本,要 readable
\begin{lstlisting}[style=jsonstyle]
{
  "target": "W",
  "location": {"lon": 121.4, "lat": -28.3},
  "g_score": 0.78,
  "tier": "High",
  "tier_reason": "score in top decile of fitted background",
  "top_experts": [
    {"name": "w_pathfinder",
     "score": 0.82, "tree_weight": 0.35,
     "effective_contrib": 0.287}
  ],
  "top_signals": [
    {"feature": "ratio_W_Mo", "z_score": 2.31,
     "source": "sediment", "expert": "w_pathfinder"},
    {"feature": "Bi", "z_score": 1.87,
     "source": "sediment", "expert": "w_pathfinder"}
  ],
  "audit": {
    "active_sources": ["wamex_sediment",
                       "wa_magnetics_grid",
                       "fault_density"],
    "active_experts": ["w_pathfinder",
                       "w_enrichment",
                       "magnetic_low",
                       "fault_proximity"],
    "model_seed": 42
  }
}
\end{lstlisting}

% 中文注释: system prompt 子节
\subsection{Faithfulness System Prompt}
\label{sec:prompt}

% 中文注释: 贴 5-6 行核心约束
The system prompt fixes the agent's behaviour with respect to numeric
and mechanistic grounding:

\begin{lstlisting}[style=jsonstyle]
You are a mineral exploration assistant for
Western Australia. You may answer questions ONLY
using values returned by tools in this turn.

HARD RULES:
1. Never fabricate scores, z-scores, coordinates,
   or active-expert lists.
2. Every numeric claim must trace verbatim to a
   tool return in this turn.
3. Never invent geological mechanisms not
   supported by reported feature signals.
4. When explaining, always cite active_experts
   and top_signals from the audit field.
5. If asked to explain, call explain_point or
   explain_region; do NOT rely on score_point
   alone.
\end{lstlisting}

% =====================================================================
% §5 Case Study Library and Demo Scenarios (~3/4 page)
% =====================================================================
\section{Case Study Library and Demo Scenarios}
\label{sec:casestudy}

% 中文注释: 这一节兼具 "数据资源贡献" 和 "demo 形态展示"

\subsection{The 100-Case Library}

% 中文注释: 9 region x 11 配置 ≈ 100,讲覆盖逻辑
We pre-compute a library of approximately 100 case studies covering
five major cratons (Yilgarn, Pilbara, Capricorn, Albany-Fraser,
Kimberley) and four well-known mining regions (Eastern Goldfields,
Murchison, Pilbara Iron, Telfer). Each named region admits two query
modes: per-commodity assessment (9 commodities $\times$ 9 regions =
81 case studies) and multi-commodity ranking (9 regions = 9 case
studies). An additional 10 cases are curated to support
specific demo dialogues (analog search anchored on Telfer; comparison
across cratons; targeted undrilled-anomaly queries).
% 中文注释: 留个 bullet 点解释 region polygon 怎么来的 - 用 CRATON 字段反推凸包
Regional polygons are approximated by convex hulls of WAMEX sediment
points carrying the corresponding \texttt{CRATON} attribute, which
avoids dependency on external shapefile licensing.

% 中文注释: 为什么预计算
\paragraph{Why precompute?}
Live regional grid scoring at the 30~$\times$~30 coarse to
100~$\times$~100 fine resolution we use elsewhere takes 8--15
seconds per (region, commodity) pair, which is acceptable in batch
contexts but disrupts conversational flow at demo time. Precomputed
case studies also permit deeper analyses---longer grids, larger
top-$K$ lists, more thorough regional narratives---that would be
prohibitive in interactive contexts.

% 中文注释: 也提一下 case study 的 schema (简短)
Each case study is a single JSON file of approximately 40~KB
containing the region's geometry, a score distribution summary
(quartiles, $p_{90}$, $p_{99}$), the top-50 anomalies with full
\texttt{explain\_point} payloads, the regional dominant pathfinders,
and a curated geological context narrative.

\subsection{Live Region Support}

% 中文注释: 自定义 region 走 live path
User-specified bounding-box and circular regions trigger live
two-stage grid scoring: a coarse 30~$\times$~30 sweep identifies the
top 5\% by score, then a 100~$\times$~100 fine grid is rasterised
within the top-region neighborhood. The two-stage scheme keeps
median latency for a $\sim$100~km radius circle under 12 seconds on a
single CPU thread.

% 中文注释: 4 个 demo scenario - Figure 3 拼图
\subsection{Demonstration Scenarios}
\label{sec:demos}

\begin{figure*}[t]
  % 中文注释: Figure 3 - 4 个对话场景拼图,每个标注关键 tool call 序列
  % 中文注释: 这是论文的视觉中心,审稿人会停在这里
  \centering
  \includegraphics[width=0.95\textwidth]{example-image-b}
  \caption{\textbf{Four representative \textsc{GroundedGeo}
  dialogues.} Each panel shows a user query (top), the agent's
  abbreviated tool-call trace (middle), and the rendered response
  (bottom). Numeric content is highlighted; arrows indicate the
  audit-field origin of each highlighted value. \textbf{(a)}~Regional
  reconnaissance over the Yilgarn Craton.
  \textbf{(b)}~Tenement-scale evaluation of a custom bounding box.
  \textbf{(c)}~Targeted undrilled-anomaly query.
  \textbf{(d)}~Why-question on regional Au prospectivity in the
  Eastern Goldfields.}
  \label{fig:scenarios}
\end{figure*}

We highlight four scenarios chosen to span the realistic working
range of an exploration geologist
(Figure~\ref{fig:scenarios}).
% 中文注释: 每个 scenario 一两句话,呼应 figure
% 中文注释: 这是论文里 "use case" 类内容,审稿人会期待看到 query/response 实例
\textbf{Regional reconnaissance} surveys all nine commodities over a
craton-scale named region; the agent invokes
\texttt{rank\_targets\_in\_region} once, then drills into the
top-ranked commodity via \texttt{top\_anomalies}.
\textbf{Tenement-scale evaluation} accepts a user-drawn bounding box
and returns a ranked list of drill targets with per-point
explanations.
\textbf{Targeted undrilled-anomaly query} layers a spatial filter
(no known site within 5\,km) onto \texttt{top\_anomalies}, surfacing
candidates not yet recorded in WAMEX.
\textbf{Why-question} prompts \texttt{explain\_region} to verbalise
the dominant pathfinder profile, distinguishing geochemical from
geophysical from structural drivers within the region.

% 中文注释: 一个具体对话节选(短)放正文,完整的进 appendix
\paragraph{Dialogue excerpt.}
A condensed exchange from scenario (a):
% 中文注释: 注意这里要用 \\ 而不是真换行,避免 LaTeX 解析问题
\begin{quote}\small\itshape
\textbf{User:} Which metals look most prospective in the Yilgarn,
and where are the hotspots?\\
\textbf{Agent:} \texttt{[calls rank\_targets\_in\_region(region=
"Yilgarn Craton")]} Ni (g-score $p_{90}=0.91$) and Au
($p_{90}=0.86$) are the top-ranked commodities. The top three
Ni anomalies are clustered around the Kalgoorlie--Norseman belt
(see top\_anomalies for full list)\ldots
\end{quote}

% =====================================================================
% §6 Faithfulness Evaluation (~1 page)
% =====================================================================
\section{Faithfulness Evaluation}
\label{sec:eval}

% 中文注释: 这一节是论文的实验,审稿人最看重的部分

\subsection{Protocol}

% 中文注释: 30 个 prompt - 来源说明
\paragraph{Prompt set.}
We curate 30 prompts authored by two geology graduate students with
field experience in Western Australia, covering point queries
($n=8$), regional reconnaissance ($n=8$), targeted undrilled-anomaly
queries ($n=8$), and analog / comparison queries ($n=6$). Prompts
intentionally include both well-known (e.g., Yilgarn Au) and
under-explored (e.g., Capricorn Ta) targets to vary the difficulty of
recalling correct geological context.

% 中文注释: 两个 condition - prompt only vs ours
\paragraph{Conditions.}
% 中文注释: 注意 baseline 要给同样的坐标 / region 信息,这是公平对比的前提
Every prompt is answered under two conditions.
(A)~\textbf{Prompt-only baseline.} GPT-4o (\texttt{gpt-4o-2024-08-06})
receives the same prompt with no tools but with the same coordinate,
region name, and target commodity information that would otherwise be
provided to the agent.
(B)~\textbf{\textsc{GroundedGeo}.} The same base model receives the
full agent loop with all twelve tools and the faithfulness system
prompt.

\subsection{Metrics}

% 中文注释: 三个 metric, 每个一段
\paragraph{Numeric hallucination rate.}
% 中文注释: 这是论文的杀手 metric
We extract every numeric claim (g-score, z-score, percentage, count,
distance) from each response using a regex over digit patterns. For
each claim, an independent verification script (i)~looks for that
number in the corresponding tool returns of the same turn (condition
B) or in the underlying scored record (condition A), and (ii)~allows
a 5\% tolerance to accommodate rounding. A claim that finds no match
is counted as a hallucination. We report mean hallucinations per
response.

\paragraph{Mechanism plausibility.}
% 中文注释: 人评 - 3 个 geology student blind rating
Three graduate students in mineral exploration, blinded to
condition, rate each response on a 1--5 Likert scale for
\emph{mechanistic plausibility} (``Are the asserted geological
processes consistent with what is known about this commodity in this
region?''). Inter-rater agreement is reported as Krippendorff's
$\alpha$.

\paragraph{Tool-call accuracy.}
% 中文注释: 只对 condition B
For condition B only, we annotate each prompt with the set of tools a
correct response \emph{should} invoke (e.g., a regional reconnaissance
prompt should invoke \texttt{rank\_targets\_in\_region} at minimum).
Tool-call accuracy is the fraction of prompts whose agent traces are
a superset of the annotated minimum.

\subsection{Results}

% 中文注释: Table 2 - 主结果
% 中文注释: 占位数字,跑完真实评估再填
\begin{table}[t]
  \centering
  \footnotesize
  \begin{tabular}{lcc}
    \toprule
    \textbf{Metric} & \textbf{Prompt-only} & \textbf{Ours} \\
    \midrule
    Numeric hallucinations / response & XX.X & \textbf{X.X} \\
    Mechanism plausibility (1--5) & X.XX & \textbf{X.XX} \\
    Tool-call accuracy & --- & \textbf{XX.X\%} \\
    Pairwise human preference & --- & \textbf{XX/30} \\
    Krippendorff $\alpha$ (plausibility) & \multicolumn{2}{c}{X.XX} \\
    \bottomrule
  \end{tabular}
  \caption{Faithfulness evaluation results on the 30-prompt
  benchmark. Numeric hallucinations are claims not traceable to
  scored data within 5\%~tolerance. Results to be populated after
  full evaluation run.}
  \label{tab:eval}
\end{table}

% 中文注释: 一段总结 main result
% 中文注释: 把数字填进来之后写得越具体越好,不要套话
Table~\ref{tab:eval} summarises results. Across the 30 prompts,
prompt-only GPT-4o produces on average \textbf{XX.X} unverifiable
numeric claims per response, dropping to \textbf{X.X} under
\textsc{GroundedGeo}---a relative reduction of \textbf{XX\%}.
Mechanism plausibility increases from \textbf{X.XX} to \textbf{X.XX}
on the 1--5 scale ($p < 0.01$, Wilcoxon signed-rank). Tool-call
accuracy is \textbf{XX\%}, with the most frequent error being
omission of \texttt{explain\_region} on why-questions.

% 中文注释: Figure 4 - side-by-side 对话, 标红幻觉
\begin{figure}[t]
  \centering
  \includegraphics[width=\columnwidth]{example-image-c}
  \caption{\textbf{Representative paired responses.} Left:
  prompt-only GPT-4o output, red overlay marks unverifiable numeric
  and mechanistic claims. Right: \textsc{GroundedGeo} response to the
  same prompt with green overlay marking the audit-field origin of
  each quantitative claim.}
  \label{fig:eval-sxs}
\end{figure}

% 中文注释: ablation - 必备子节,即便只一小段
\subsection{Ablations}

% 中文注释: 3 个 ablation 设计, 每个一句话
We ablate three design choices on a subset of 12 prompts.
(i)~\textbf{Removing audit fields} from tool returns increases
numeric hallucinations from \textbf{X.X} to \textbf{X.X}, confirming
that agents lean on \texttt{audit} for quotation rather than only the
top-level fields.
(ii)~\textbf{Removing the faithfulness system prompt} (keeping
tools) increases hallucinations to \textbf{X.X} while leaving
tool-call accuracy nearly unchanged.
(iii)~\textbf{Replacing case-study returns with raw grid scores}
increases response latency by 8$\times$ but does not affect
faithfulness.

% =====================================================================
% §7 Discussion (~1/4 page)
% =====================================================================
\section{Discussion}
\label{sec:discussion}

% 中文注释: 这一节短,主要给两个 takeaway

\paragraph{Where faithfulness comes from.}
% 中文注释: 把功劳归在 architecture 而不是 prompt
The structural absence of numeric hallucination in \textsc{GroundedGeo}
is not primarily attributable to prompt engineering: ablating the
faithfulness prompt while keeping audit-rich tools leaves
hallucinations rare. Rather, hallucination is suppressed because the
LLM has no in-context source of unverified numbers---every digit it
might quote is already tagged with its origin in some tool return.
This suggests a generalisable design pattern: \emph{tool returns
should be the only quantitative information surface available to the
agent}, with prompts playing a secondary reinforcing role.

\paragraph{What the LLM still contributes.}
% 中文注释: 防御 'LLM 没用' 的可能审稿评论
A natural concern is that \textsc{GroundedGeo} simply uses the LLM as
a thin verbalisation wrapper. We observe three non-trivial
contributions: (1)~tool composition planning (interleaving
\texttt{rank\_targets\_in\_region} with multiple \texttt{top\_anomalies}
calls); (2)~surface realisation that compresses 12--15 tool returns
into a 4-sentence response; (3)~handling ambiguous user references
(``the gold region around Kalgoorlie'' resolves to a circular region
spec). Removing the LLM and substituting a template would lose all
three capabilities.

% =====================================================================
% §8 Limitations + Ethics
% =====================================================================
\section*{Limitations}

% 中文注释: ACL 要求必备 Limitations section
% 中文注释: 写三类: 数据局限, 方法局限, 评估局限

\paragraph{Geographic scope.}
The case study library covers only Western Australia. Extending to
other jurisdictions requires re-running the scoring pipeline on local
data sources and re-fitting the pathfinder dictionaries; we estimate
roughly 2--3 person-weeks per jurisdiction.

\paragraph{Template-bounded narrative diversity.}
Surface realisations in \texttt{explain\_point} and
\texttt{explain\_region} are template-based, which preserves
faithfulness but limits stylistic variation. The agent introduces
some diversity through paraphrase, but adversarial prompts requesting
elaborate explanations may surface template repetition.

\paragraph{Evaluation prompt distribution.}
Our 30-prompt benchmark is authored by two annotators and may not
fully represent the distribution of real exploration queries. We
release the prompt set publicly to enable benchmark extension.

\paragraph{Not for commercial decisions.}
% 中文注释: 这是地学领域专用的免责声明 -- 审稿人会赞赏看到
\textsc{GroundedGeo} produces statistical prospectivity scores. It
should not be interpreted as a substitute for ground-truth drilling,
sampling, or licensed mineral resource estimation procedures.

\section*{Ethics Statement}

% 中文注释: ACL 要求 Ethics statement(可短)
% 中文注释: 主要讲: 数据 license, 不涉及人类受试者, 但提到 dual-use 风险
All data used in \textsc{GroundedGeo} are publicly available under
the WAMEX open data license. No human-subject data are involved.
We note that high-resolution prospectivity maps can in principle
influence mineral tenement competition; we believe the benefits to
research transparency outweigh this risk for an open data
jurisdiction, but caution against direct deployment in jurisdictions
with restricted geoscientific data.

\section{Conclusion}

% 中文注释: 收尾段,1 段
We presented \textsc{GroundedGeo}, an LLM agent for mineral
prospectivity reasoning where every quantitative and mechanistic
claim is structurally traceable to a typed tool return. By combining
audit-rich tools, a precomputed case study library spanning nine
commodities and nine regions of Western Australia, and a faithfulness
system prompt, the system reduces numeric hallucinations by \textbf{XX\%}
relative to prompt-only GPT-4o while preserving fluent natural-language
interaction. We hope the pattern of \emph{quantitative-information-
surface restriction via tool typing} generalises to other scientific
domains where the underlying data are accessible but free-form LLM
outputs are not yet trusted.

% =====================================================================
% References
% =====================================================================
% 中文注释: refs 要写真实存在的引文,以下是要写进 custom.bib 的关键文献清单
% - yao2022react              : Yao et al. ReAct (ICLR 2023)
% - schick2023toolformer      : Schick et al. Toolformer (NeurIPS 2023)
% - openai2023functioncalling : OpenAI function calling release blog/doc
% - mialon2023gaia            : Mialon et al. GAIA benchmark
% - liu2023agentbench         : Liu et al. AgentBench
% - reiter2000building        : Reiter & Dale, Building NLG Systems
% - lewis2020retrieval        : Lewis et al. RAG
% - bohnet2022attributed      : Bohnet et al. Attributed QA
% - min2023factscore          : Min et al. FActScore
% - singhal2023large          : Singhal et al. Med-PaLM
% - dahl2024large             : Dahl et al. legal hallucinations
% - petroni2021kilt           : Petroni et al. KILT
% - carranza2017improved      : Carranza, improved WofE
% - zuo2020deep               : Zuo, deep learning for prospectivity
% - sun2020data               : Sun et al. data-driven prospectivity
% - wang2022deep              : Wang et al. SHAP on RF prospectivity

\bibliography{custom}

% =====================================================================
% Appendix
% =====================================================================
\appendix

\section{Full Tool Schemas}
\label{sec:appendix-tools}

% 中文注释: 12 个工具的完整 OpenAI function calling JSON schema 进 appendix
% 中文注释: 这是 reproducibility 要求

Full OpenAI function-calling JSON schemas for all twelve tools are
included in the supplementary archive. Below we reproduce the schema
for \texttt{top\_anomalies} as a representative example.

\begin{lstlisting}[style=jsonstyle]
{
  "name": "top_anomalies",
  "description": "Return top-K prospectivity
   anomalies within a region for a target metal,
   optionally excluding points near known sites.",
  "parameters": {
    "type": "object",
    "required": ["region", "target"],
    "properties": {
      "region": {"$ref": "#/Region"},
      "target": {"type": "string",
                 "enum": ["Cu","Au","Ni","W",
                          "Sn","Co","Ta","Mn","REE"]},
      "top_k": {"type": "integer", "default": 10},
      "exclude_known": {"type": "boolean",
                        "default": true},
      "min_distance_km": {"type": "number",
                          "default": 5.0}
    }
  }
}
\end{lstlisting}

\section{Evaluation Prompt Set}
\label{sec:appendix-prompts}

% 中文注释: 30 prompts 全文 - reproducibility 要求
The full list of 30 evaluation prompts is included in the
supplementary archive. We reproduce two representative prompts:

\begin{quote}\itshape
\textbf{P07 (regional reconnaissance):} ``Looking at the Capricorn
Orogen, which of the nine target commodities show the strongest
prospectivity, and what are the dominant geochemical signatures?''

\textbf{P19 (undrilled anomaly):} ``Find tungsten anomalies in the
Eastern Goldfields with a g-score above 0.7 that have no known
W deposit within 10 km.''
\end{quote}

\section{Human Evaluation Interface}
\label{sec:appendix-human}

% 中文注释: 简短描述 human eval 是怎么做的 - 给 inter-rater agreement 提供 context
Annotators viewed paired responses in a randomised, condition-blinded
side-by-side interface and rated each on the 1--5 plausibility scale
before indicating a pairwise preference. Each prompt was rated by all
three annotators; we report mean ratings and Krippendorff's $\alpha$
on the ordinal scale.

\end{document}
